%!TEX root = ../../note.tex
% Source: ../note-course/chapters/u2/s4.tex
\subsection{Suprema and infima}
Bounds make sense in any ordered set. What is special about $\R$ is that
every nonempty set that is bounded above has a supremum in $\R$; this need
not hold in $\Q$.

\begin{defi}[Bounded sets]
  Let $S$ be a nonempty subset of $\R$. Then $S$ is
  \begin{enumerate}
    \item \term{bounded above} if there is $u \in \R$ (an \term{upper bound})
      with $x \leq u$ for all $x \in S$;
    \item \term{bounded below} if there is $w \in \R$ (a \term{lower bound})
      with $w \leq x$ for all $x \in S$;
    \item \term{bounded} if it is bounded above and below, and
      \term{unbounded} otherwise.
  \end{enumerate}
\end{defi}

\begin{defi}[Supremum and infimum]
  An upper bound $u_0$ of $S$ such that $u_0 \leq u$ for every upper bound
  $u$ of $S$ is the \term{supremum} (or \emph{least upper bound}) of $S$,
  written $u_0 = \sup S$.

  A lower bound $w_0$ of $S$ such that $w_0 \geq w$ for every lower bound $w$
  of $S$ is the \term{infimum} (or \emph{greatest lower bound}) of $S$,
  written $w_0 = \inf S$.
\end{defi}

The notation is justified because a set has at most one supremum: if $u_1$
and $u_2$ are both suprema of $S$, then $u_1 \leq u_2$, since $u_1$ is least
and $u_2$ is an upper bound, and symmetrically $u_2 \leq u_1$. The same
argument with the inequalities reversed shows that the infimum is unique.

\subsubsection*{The \texorpdfstring{$\epsilon$}{epsilon}-characterization of the supremum}
An upper bound $u$ fails to be the \emph{least} one exactly when some smaller
number is still an upper bound, i.e.\ when there is a gap between $S$ and
$u$. So $u$ is the supremum precisely when there is no gap: however small a
margin $\epsilon > 0$ we allow, $S$ has a point within $\epsilon$ of $u$.

\begin{nlemma}[$\epsilon$-characterization of the supremum]\label{lem:eps-sup}
  An upper bound $u$ of a nonempty set $S \subseteq \R$ is the supremum of
  $S$ if and only if for every $\epsilon > 0$ there is $x_\epsilon \in S$
  with $u - \epsilon < x_\epsilon$.
\end{nlemma}

\begin{proof}
  ($\Rightarrow$) Let $u = \sup S$ and $\epsilon > 0$. Since
  $u - \epsilon < u$ and $u$ is the \emph{least} upper bound, $u - \epsilon$
  is not an upper bound. So some $x_\epsilon \in S$ satisfies
  $u - \epsilon < x_\epsilon \leq u$.

  ($\Leftarrow$) Suppose $u$ is not least, so there is an upper bound
  $w < u$. Take $\epsilon = u - w > 0$. Then some $x_\epsilon \in S$
  satisfies $w = u - \epsilon < x_\epsilon$, so $w$ is not an upper bound
  after all.
\end{proof}

So to show $u = \sup S$, one checks that $u$ is an upper bound of $S$, and
that for every $\epsilon > 0$ there is $x \in S$ with $u - \epsilon < x \leq u$.
For $\inf S$, reverse the inequalities: a lower bound $w$ such that every
$\epsilon > 0$ has some $x \in S$ with $w \leq x < w + \epsilon$.

\begin{eg}
  Let $S = \{x \in \R : 0 \leq x < 1\}$. Then $\sup S = 1$ and $\inf S = 0$.

  \emph{$\sup S = 1$}: every $x \in S$ has $x < 1$, so $1$ is an upper bound.
  Since $S$ is nonempty and bounded above, $u = \sup S$ exists by
  completeness, and $u \leq 1$. Suppose $u < 1$, and let
  $\epsilon = \frac{1 - u}{2} > 0$. Then
  \[
    u < u + \epsilon = \frac{u + 1}{2} < 1,
  \]
  and $u + \epsilon \geq 0$ since $u \geq 0$ (as $0 \in S$). So
  $u + \epsilon \in S$ exceeds the upper bound $u$, a contradiction. Hence
  $u = 1$.

  \emph{$\inf S = 0$}: every $x \in S$ has $x \geq 0$, so $0$ is a lower
  bound, and $w = \inf S \geq 0$. Every lower bound satisfies $w \leq 0$
  because $0 \in S$. Hence $w = 0$.
\end{eg}

\begin{warning}
  $\sup S$ need not belong to $S$. Above, $\sup S = 1 \notin S$, so $S$ has
  no maximum; $\inf S = 0 \in S$, so $0$ is also the minimum. A supremum is a
  maximum only when it belongs to the set. Note also that ``bounded above'',
  ``bounded'' and so on are properties of the \emph{set} $S$, whereas
  $\sup S$ and $\inf S$ (when they exist) are \emph{real numbers}.
\end{warning}
